% Options for packages loaded elsewhere
\PassOptionsToPackage{unicode}{hyperref}
\PassOptionsToPackage{hyphens}{url}
%
\documentclass[
  10pt,
]{article}
\usepackage{amsmath,amssymb}
\usepackage{iftex}
\ifPDFTeX
  \usepackage[T1]{fontenc}
  \usepackage[utf8]{inputenc}
  \usepackage{textcomp} % provide euro and other symbols
\else % if luatex or xetex
  \usepackage{unicode-math} % this also loads fontspec
  \defaultfontfeatures{Scale=MatchLowercase}
  \defaultfontfeatures[\rmfamily]{Ligatures=TeX,Scale=1}
\fi
\usepackage{lmodern}
\ifPDFTeX\else
  % xetex/luatex font selection
\fi
% Use upquote if available, for straight quotes in verbatim environments
\IfFileExists{upquote.sty}{\usepackage{upquote}}{}
\IfFileExists{microtype.sty}{% use microtype if available
  \usepackage[]{microtype}
  \UseMicrotypeSet[protrusion]{basicmath} % disable protrusion for tt fonts
}{}
\makeatletter
\@ifundefined{KOMAClassName}{% if non-KOMA class
  \IfFileExists{parskip.sty}{%
    \usepackage{parskip}
  }{% else
    \setlength{\parindent}{0pt}
    \setlength{\parskip}{6pt plus 2pt minus 1pt}}
}{% if KOMA class
  \KOMAoptions{parskip=half}}
\makeatother
\usepackage{xcolor}
\usepackage[margin=0.85in]{geometry}
\setlength{\emergencystretch}{3em} % prevent overfull lines
\providecommand{\tightlist}{%
  \setlength{\itemsep}{0pt}\setlength{\parskip}{0pt}}
\setcounter{secnumdepth}{-\maxdimen} % remove section numbering
\usepackage{amsmath,amssymb,amsthm}
\usepackage{longtable}
\ifLuaTeX
  \usepackage{selnolig}  % disable illegal ligatures
\fi
\usepackage{bookmark}
\IfFileExists{xurl.sty}{\usepackage{xurl}}{} % add URL line breaks if available
\urlstyle{same}
\hypersetup{
  pdftitle={Strict tropical barriers and conditional singular Runge reconstruction},
  pdfauthor={Primary-source reconstruction; unverified mathematical draft},
  hidelinks,
  pdfcreator={LaTeX via pandoc}}

\title{Strict tropical barriers and conditional singular Runge
reconstruction}
\author{Primary-source reconstruction; unverified mathematical draft}
\date{5 October 2026}

\begin{document}
\maketitle

\subsection{1. Disposition and the unchanged
problem}\label{disposition-and-the-unchanged-problem}

The original P16.1 is not solved by this paper. The candidate flag is
false. The original problem is:

\begin{quote}
P16.1 (T5): Finite-energy equations on singular or unibranch varieties
Extend the finite-energy non-Archimedean Monge--Ampère operator and its
existence-and- uniqueness theorem to singular or unibranch varieties,
replacing each use of smoothness in compactness, envelopes, energy
differentiation and comparison. State all retained polarization and field
hypotheses.
\end{quote}

Its unchanged completion criterion is:

\begin{quote}
Establish the genuine finite-energy operator, exact measure mass and
energy condition, existence, uniqueness modulo constants, and every
compactness/envelope/differentiation/comparison replacement for the
singular/geometrically unibranch projective domain with explicitly
retained field and polarization hypotheses. Verify the new envelope
source and the actual general-variety global theory. No change to the
original mathematical target is permitted.
\end{quote}

The full family, task object and assignment are preserved in
\texttt{input.json}, its byte-identical copy
\texttt{input-original.json}, \texttt{task-original.json} and
\texttt{assignment-original.txt}. The equation domain retained here is
every reduced irreducible geometrically unibranch projective variety
\(X\) over every algebraically closed trivially valued
characteristic-zero field \(k\), with every ample line bundle \(L\),
every dimension, and arbitrary-support targets satisfying the intrinsic
energy condition below. Auxiliary countability and rationality
hypotheses in this paper do not narrow that domain.

There are three distinct kinds of results. First, we verify the actual
scope and use of the established global theory in Boucksom--Jonsson
(BJ). Its finite-energy operator, comparison and uniqueness already have
general-variety scope. Existence is conditional on envelopes. Second, we
prove explicit strict barriers, real-support approximations,
original-ring pole-order localization, strict Hartogs radius formulae
and a finite-projection lemma. Third, we reconstruct a repaired
split-piece Runge implication under explicitly stated metric, smooth
regularization and trace antecedents. Those antecedents are not imported
from generated rounds as theorems. Their complete independent
reconstruction, the countable singular envelope theorem and hence
unconditional singular existence remain missing.

The embedded round70 paper and notes, round79 proposal, and round71
diagnostic have all been read as supplied mathematical data. They are
retained locally in \texttt{embedded-data}. No filesystem path printed
in those documents has been accessed. The ten authorized primary caches
have exact copies and provenance. The Wu preprint is version one, dated
September 2026; its statements are not treated as correctness
certificates. None of the mathematical drafts in this paper has an
independent correctness verdict. Successful execution and compilation
are observations only.

\subsection{2. The actual global operator and the equation
reduction}\label{the-actual-global-operator-and-the-equation-reduction}

Write \(n=\dim X\), \(\omega=c_1(L)\) and \(V=(L^n)>0\). Let
\(\operatorname{PSH}(X,\omega)\) always mean BJ's genuine global
generated cone, not the corrected auxiliary local cone introduced later.
For this section geometric unibranchness is unnecessary. Irreducibility
gives the one-constant uniqueness statement. BJ also treats ample real
numerical Cartier classes; the original line-bundle target suffices
here.

\subsubsection{2.1. Model geometry retains a nonnormal generic
fiber}\label{model-geometry-retains-a-nonnormal-generic-fiber}

BJ v4 LF450--469 defines varieties to be reduced, with no normality
assumption. Section 7 explicitly resumes general projective varieties at
LF3649--3658. This scope is decisive: there is a genuine singular-space
operator already in the source.

On a test configuration \(\mathcal X\to\mathbb A^1\), integral closure
means closure \textbf{inside its prescribed generic fiber}. On an affine
chart \(A\), close \(A\) in \(A\otimes_{k[t]}k(t)\). BJ Definition 1.10,
Example 1.11 and Theorem 1.13, LF759--849, prove finiteness and the
vertical \(R_1,S_2\) criterion. Absolute normalization would change a
nonnormal generic fiber and is not substituted.

Here is the relevant mechanism in the proof. A vertical fractional ideal
has poles only on the central fiber. At an associated polar point, an
element \(f\) satisfies \(f\mathfrak m\mathfrak a\subset A\) but
\(f\mathfrak a\not\subset A\). If the former ideal lies in
\(\mathfrak m\), the determinant trick would make \(f\mathfrak a\)
integral and hence regular, a contradiction. Otherwise
\(f\mathfrak m\mathfrak a=A\), making \(\mathfrak m\) invertible, so the
point is a DVR. Conversely vertical depth at least two rules out a
higher-codimension associated pole, while the vertical DVRs rule out
codimension-one poles. This is a replacement for absolute normality, not
a smoothness assumption.

For a central component \(E\), let \(b_E=\operatorname{ord}_E(t)>0\) and
\(v_E=b_E^{-1}\operatorname{ord}_E|_{k(X)}\). With relatively nef model
data \(\theta_{i,\mathcal X}+D_i\), the actual raw mixed measure is \[
 M^{\mathrm{raw}}((\theta_1,\phi_1),\ldots,(\theta_n,\phi_n))
 =\sum_E b_E\prod_{i=1}^n((\theta_{i,\mathcal X}+D_i)|_E)\,\delta_{v_E}. \tag{2.1}
\] The product denotes an intersection number, including the central
fundamental cycle. Nefness gives nonnegative weights. Rational
equivalence of the central and generic fibers gives total mass
\((\theta_1\cdots\theta_n)_X\). In particular the diagonal raw mass is
\(V\), not one. The factor \(b_E\) cannot be suppressed.

\subsubsection{2.2. Energy and Radon
representation}\label{energy-and-radon-representation}

The compactified intersection pairing is symmetric and independent of
determination by the projection formula. Translating its \(i\)th
potential by \(c\) adds \(c(\prod_{j\ne i}\theta_j)_X\). BJ 7.1 extends
it as the infimum over model majorants. Monotonicity follows from nef
intersections. To verify upper semicontinuity, choose a majorant tuple
below a prescribed strict pairing threshold. The functionals
\(\sup(\phi_i-\psi_i)\) are continuous in BJ's divisorial topology, so
small translated majorants remain above all nearby potentials. The
translation rule controls their pairing. This proves upper
semicontinuity and decreasing-net continuity; it uses no envelope
property.

BJ 7.8--7.14 extend the finite mixed pairing by multilinearity to
differences of finite-energy potentials. Ample class comparisons and the
expansion of the energy of sums ensure that the requisite mixed terms
are finite. Define \[
 E_\omega(\phi)=\frac{\langle(\omega,\phi)^{n+1}\rangle}{(n+1)V},\qquad
 \mathcal E^1(X,\omega)=\{\phi\in\operatorname{PSH}(X,\omega):E_\omega(\phi)>-\infty\}.
 \tag{2.2}
\] We have \(E(0)=0\), \(E(\phi+c)=E(\phi)+c\), and \(E\) is increasing,
concave, upper semicontinuous and decreasing-continuous.

For a genuine algebraic PL test \(a\), set \[
 \int a\,dM^{\mathrm{raw}}
 =\langle(0,a),(\theta_1,\phi_1),\ldots,(\theta_n,\phi_n)\rangle. \tag{2.3}
\] Any PL test is a difference of positive model potentials in an
auxiliary ample class. Positivity bounds the functional by its mass
times \(\|a\|_\infty\). Uniform PL density and Radon representation give
the genuine measure. This is precisely the construction and proof in BJ
7.18--7.19, LF3980--4057. It gives symmetry, multilinearity, model
agreement, exact mass, mixed integrability and decreasing weak
continuity. Define the probability convention
\[\operatorname{MA}_\omega(\phi)=V^{-1}M^{\mathrm{raw}}_\omega(\phi). \tag{2.4}\]

Weighted continuity requires a sign repair in the printed proof of 7.21:
the pairing containing the test is subtracted \textbf{from neither}
integral; rather \[
 \int \phi_0\,dM^{\mathrm{raw}}
 =\langle(\omega_0,\phi_0),\ldots\rangle
  -\langle(\omega_0,0),\ldots\rangle. \tag{2.5}
\] This is the identity already printed correctly at LF4107--4113. The
opposite order at LF4119--4123 is inconsistent with (2.3). Both
corrected pairings converge. Translate \(\phi_0\) below zero and apply
the positive weighted-measure lemma to \(-\phi_0\), obtaining weighted
weak convergence. This repair does not assume the measure is monotone in
the potential.

\subsubsection{2.3. Comparison, domination and uniqueness precede
EP}\label{comparison-domination-and-uniqueness-precede-ep}

On a strict model winning component, maxima and the winning potential
have the same divisor near that component, hence the same intersection
weight. To pass to finite energy, write
\[ f_j=\max(\phi_j,\psi_j)-\psi_j.\] Apply weighted continuity
separately to the two decreasing psh terms. The difference \(f_j\) need
not decrease. This supplies the missing weighted argument behind BJ
LF4581--4585 and gives \[
 \mathbf1_{\{\phi>\psi\}}\operatorname{MA}(\max(\phi,\psi))
 =\mathbf1_{\{\phi>\psi\}}\operatorname{MA}(\phi). \tag{2.6}
\] Equal masses and the displaced maximum \(\max(\phi,\psi-\epsilon)\)
give \[\int_{\{\phi<\psi\}}\operatorname{MA}(\psi)
 \le\int_{\{\phi<\psi\}}\operatorname{MA}(\phi). \tag{2.7}\]
Ordinary-open mixed locality follows by displacing either potential,
decreasing limits and polarization.

Put
\(I(\phi,\psi)=\int(\phi-\psi)(\operatorname{MA}(\psi)-\operatorname{MA}(\phi))\).
Mixed intersection negativity gives \(I\ge0\) and the tangent energy
inequalities. The separation estimate BJ 10.3, applied to divisorial
finite-energy Diracs and the trivial Dirac, makes \(I=0\) imply equality
of all divisorial evaluations after equal-sup normalization. Divisorial
determination then gives equality everywhere. This is the proof of
actual \textbf{Corollary} 10.4, LF5444--5456. Equal measures therefore
imply one additive real constant on integral \(X\).

Actual \textbf{Corollary} 10.6, LF5468--5477, proves domination without
EP. Replace \(\phi\) by \(\max(\phi,\psi)\) if \(\phi\le\psi\)
MA\((\psi)\)-almost everywhere. The replacement is finite energy and
equals \(\psi\) almost everywhere for that measure; its \(I\) with
\(\psi\) is nonpositive. Separation forces a constant difference, and
the almost-everywhere equality removes the constant.

\subsubsection{2.4. Exact targets and actual attainment under
EP}\label{exact-targets-and-actual-attainment-under-ep}

A raw target is a positive Radon measure \(\nu\) on \(X^{\mathrm{an}}\)
with \(\nu(X^{\mathrm{an}})=V\). Put \(\mu=\nu/V\). Its intrinsic
finite-dual-energy condition is \[
 E_\omega^\vee(\mu)=\sup_{\psi\in\mathcal E^1(X,\omega)}
 \left(E_\omega(\psi)-\int\psi\,d\mu\right)<\infty. \tag{2.8}
\] The integral can be extended-valued in the definition; finiteness
forces every finite-energy potential to be integrable. Conversely BJ
9.20 proves integrability of every finite-energy potential implies (2.8)
for irreducible \(X\). To see the operative proof, normalize
\(\sup\psi=0\). If negative integration were unbounded on
\(J(\psi)=-E(\psi)\le C\), choose \(\psi_j\) with integrals at most
\(-2^j\). The decreasing convex sums \(\sum_{j\le m}2^{-j}\psi_j\) (with
the unused weight on zero) have energy at least \(-C\). Their limit is a
genuine finite-energy potential with integral \(-\infty\), a
contradiction. Fixed-ray mixed expansion gives
\(J(t\psi)\le C_n t^2J(\psi)\), so scaling proves
\(-\int\psi\,d\mu\le A\sqrt{J(\psi)}+B\). Taking the dual supremum
proves finiteness. There is no support restriction. Actual MA targets
satisfy (2.8) by mixed integrability and the concave tangent inequality.

BJ Theorem 8.5 is an unconditional \textbf{outer-energy} statement:
\[ E^\downarrow(\phi+tf)=E(\phi)+t\int f\operatorname{MA}(\phi)+O_f(t^2),
 \quad f\in\operatorname{PL}_{\mathbb R}, \tag{2.9}\] uniform in
\(\phi\in\mathcal E^1\). Its proof represents \(f\) by two ample model
potentials and bounds mixed measures by their fixed intersection
degrees. It passes from continuous inputs to decreasing finite-energy
inputs. BJ's proof cites its relative-volume inputs; those cited outside
the ten supplied caches are not independently re-proved here. For
continuous \(f\), PL approximation gives derivative bounds with
endpoints \(\int f_m\operatorname{MA}(\phi)\mp\epsilon_m\), correcting
the two-minus printed squeeze at LF4950--4954. The preliminary estimate
is \(|f-f_m|\le\epsilon_m\), not \(t\epsilon_m\) before multiplication
by \(t\). For an unbounded \(\phi\), identifying
\(E^\downarrow(\phi+tf)\) with \(E(P(\phi+tf))\) requires EP.

Assume genuine EP for \(\omega\). BJ 5.11 proves normalized-PSH
compactness by Tychonoff on bounded divisorial evaluations, genuine tail
upper envelopes, and nonnegligibility of divisorial points. A maximizing
sequence for (2.8) has bounded \(J\) by 9.15. Its weak subnet limit has
finite energy by \(J\) lower semicontinuity. Proposition 9.19 supplies
\(L^1(\mu)\) convergence on that bounded-\(J\) subnet: approximate the
target integration uniformly on bounded-\(J\) sets by integration
against maximizing model MA measures. This approximation precedes
attainment; it does not assume an inverse. Correct the printed
absolute-difference identity at LF5254 to \[|a-b|=2\max(a,b)-a-b.\]
Energy upper semicontinuity now yields an actual maximizer. Both signs
in (2.9), combined with outer duality, give equality of the genuine MA
measure and \(\mu\). The raw objective \(VE(\phi)-\int\phi\,d\nu\) is
invariant under constants exactly at raw mass \(V\).

Thus the source-based theorem is: \textbf{under EP}, every target (2.8)
has a genuine solution, unique modulo constants. Actual BJ 12.8,
LF6092--6141, equates EP, onto MA and completeness; its 12.9 is the
smooth sufficient endpoint. The source does not make singular EP
automatic. Dimension zero is immediate: \(X=\operatorname{Spec}k\),
\(V=1\), \(E(c)=c\), with the one-point equation.

\subsection{3. Exact auxiliary definitions and elementary
closures}\label{exact-auxiliary-definitions-and-elementary-closures}

For affine \(Y=\operatorname{Spec}A\) define actual rational affine FS
functions (AFS) to be \[\max(\lambda,c_i\log|f_i|+C_i)_i,\] where all
displayed constants and positive coefficients are rational and the list
of regular \(f_i\) contains algebra generators of \(A\). A generating
system for a finite-dimensional space or a list of powers alone is
insufficient. Upper growth means \(u\le H+C\) for some AFS \(H\).

A corrected continuous rational local contact function is a finite
continuous selection from finitely many rational affine expressions in
regular log norms on finite local charts. A contact support satisfies
\(b(s)\ge\ell(s)+C\) and \(b(t)=\ell(t)+C\), with the \textbf{same}
constant. At an infinite tropical coordinate, the support and every
branch used on its finite chart must have zero coefficient in that
coordinate; shrink so all remaining used functions are units. Decreasing
local limits on common neighborhoods give an auxiliary local class. A
weak input has one proper birational algebraic witness whose pullback is
locally in that class, and is finite at the generic trivial seminorm.
Proper closedness gives upper semicontinuity downstairs. These
definitions do not replace genuine global PSH.

Positive rational sums and maxima of AFS functions are AFS after
distributing their finite maxima and clearing rational exponents into
regular monomials. To retain each original generator, use its branch in
one factor and floors in all others. If the resulting monomial uses a
generator power, append a low dominated branch of the generator itself:
for \(c>0\), \(\eta\in(0,c)\) and a finite floor,
\(\eta\log|z|-M\le\max(\lambda,c\log|z|+C)\) for all real log norms when
\(M\) is sufficiently large. This also handles \(z=0\). Differences form
a rational vector lattice. Generator branches make AFS sublevels compact
in a closed affine embedding.

Restriction to a principal open introduces inverse generators. If
\(u_j\downarrow u\) is AFS on the old affine ring and \(H_f\ge0\) is an
AFS exhaustion including the inverse and all new generators, then
\(\max(u_j,2^{-j}H_f-j)\) decreases to \(u\) on that open and is actual
AFS there. Polynomial degree bounds give upper growth for its limit.
This supplies restrictions used in later localized inductions.

\subsection{4. Explicit strict barriers and real logarithmic
supports}\label{explicit-strict-barriers-and-real-logarithmic-supports}

\subsubsection{4.1. A strict support gap on extended tropical
coordinates}\label{a-strict-support-gap-on-extended-tropical-coordinates}

Let \(T_N=(\mathbb R\cup\{+\infty\})^N\). Choose an integer
\(a\ge\log(N+1)\) and set
\[\chi_N(t)=a^{-1}\log\left(1+\sum_{i=1}^N e^{-at_i}\right),\quad e^{-a\infty}=0. \tag{4.1}\]
It is continuous on \(T_N\), convex and coordinatewise decreasing on the
finite part, and \[\max(0,-t_1,\ldots,-t_N)\le\chi_N(t)
 \le\max(0,-t_1,\ldots,-t_N)+1. \tag{4.2}\] Fix \(t\). Let \(I\) be its
finite-coordinate set, \(Z=1+\sum_{i\in I}e^{-at_i}\), \(p_0=1/Z\) and
\(p_i=e^{-at_i}/Z\) for \(i\in I\). Define
\[\ell_t(s)=\chi_N(t)-\sum_{i\in I}p_i(s_i-t_i). \tag{4.3}\] At a pole
\(s_i=\infty\) for \(i\in I\), give this support value \(-\infty\). Its
other coefficients are zero. The convex support inequality gives
\(\chi_N\ge\ell_t\). Equality occurs only at \(s=t\): the Hessian in the
finite \(I\) coordinates is \(a(\operatorname{diag}(p)-pp^T)\), positive
definite since \(p_0>0\); any finite value in a formerly infinite
coordinate strictly increases (4.1).

We prove the global gap, not just strict pointwise contact. In finite
\(I\) coordinates the gap is bounded below by
\[\max(0,-s_i:i\in I)+\sum_{i\in I}p_i s_i-C_t.\] The homogeneous part
is positive on every nonzero vector: its first term is the maximum of
\(0,-s_i\), and the negative of its second term is a strict weighted
average of those values, with all weights \(p_i,p_0\) positive.
Compactness of the unit sphere gives coercivity. A bounded gap sublevel
therefore confines these coordinates to a bounded interval. The gap
bound and bounded support then give a lower bound for the other
coordinates; their upper endpoint \(\infty\) is allowed. Gap sublevels
are closed and compact in the extended product; when a used coordinate
tends to \(\infty\) the gap tends to \(+\infty\). The zero set is
exactly \(\{t\}\). In any compact gap sublevel, the complement of a
neighborhood of \(t\) has a positive minimum. Outside that sublevel the
gap is already uniformly positive. Hence, for every neighborhood \(V\)
of \(t\), there is \(\delta>0\) with
\[\chi_N(s)-\ell_t(s)\ge\delta\quad(s\notin V). \tag{4.4}\] If \(I\) is
empty, the support is zero; (4.1) directly bounds every coordinate below
on a gap sublevel, and the same argument applies. This proves the full
extended-coordinate statement without the printed Gaussian-extension
argument.

\subsubsection{4.2. Uniform multinomial AFS
approximants}\label{uniform-multinomial-afs-approximants}

Put \(z_0=1\), \(z_i=e^{-at_i}\). Expanding \((\sum_{i=0}^N z_i)^m\)
gives \(B_m=\binom{m+N}{N}\) nonnegative summands. Their maximum lies
between the sum divided by \(B_m\) and the sum. Thus
\[q_m(t)=\max_{|\beta|=m}\left\{-\sum_{i=1}^N\frac{\beta_i}{m}t_i
 +\frac{\log\binom m\beta}{am}\right\}\] satisfies
\[q_m\le\chi_N\le q_m+\frac{\log B_m}{am} \tag{4.5}\] on the entire
extended product. Terms meeting a pole have value \(-\infty\); the
\(\beta_0=m\) term is the zero floor. After composition
\(t_i=-\log|f_i|\), the regular monomial is \(\prod f_i^{\beta_i}\) and
its log coefficient is \(1/m\). If the \(f_i\) contain generators, the
terms \(\beta_i=1\), \(\beta_0=m-1\) contain those generators
themselves.

Choose \(m_j\) with the error in (4.5) at most \(\epsilon_j/4\), round
every finite intercept upward to rationals with error at most
\(\epsilon_j/4\), and add rational \(\epsilon_j\), where
\(\epsilon_{j+1}\le\epsilon_j/4\). The resulting AFS value is between
\(\chi_N+3\epsilon_j/4\) and \(\chi_N+5\epsilon_j/4\). Consequently the
sequence decreases everywhere and converges to \(\chi_N\). This is
global uniform approximation, including poles.

\subsubsection{4.3. Real supports require individual
floors}\label{real-supports-require-individual-floors}

Let \(p_i>0\) be arbitrary real coefficients, \(C\in\mathbb R\), and
\(f_i\) regular. Choose rational AFS \(H\ge1\) containing the original
generators and integers \(D_i\ge1\) with \(\log|f_i|\le D_iH\). Define
\[F_{ij}=\max(\log|f_i|,2^{-j}H-j). \tag{4.6}\] These are actual AFS,
decrease to the individual logs, and obey \(F_{ij}\le D_iH\). Choose
nested rational intervals \(0\le q_{ij}\le p_i\le r_{ij}\), with
\(q_{ij}\uparrow p_i\) and \(r_{ij}\downarrow p_i\), and rational
\(C_j\downarrow C\). Put \[\eta_j=\sum_iD_i(r_{ij}-q_{ij}),\qquad
 L_j=\sum_iq_{ij}F_{ij}+(\eta_j+2^{-j})H+C_j. \tag{4.7}\] For nonzero
\(f_i\) at the point, replacing \(p_i\) by \(q_{ij}\) costs at most
\(D_i(p_i-q_{ij})H\); negative logs only improve this inequality. Thus
\(L_j\ge\sum p_i\log|f_i|+C\). At a pole the inequality is immediate.
Monotonicity is exact:
\[\eta_j-\eta_{j+1}\ge\sum_iD_i(q_{i,j+1}-q_{ij}).\] Expand
\(L_j-L_{j+1}\), use \(F_{ij}\ge F_{i,j+1}\) and \(F_{i,j+1}\le D_iH\),
and the displayed inequality gives a nonnegative result. The positive
\(2^{-j}H\) term and the individual lower bounds in (4.6) retain
generator growth; after distribution low generator branches can be
appended as in Section 3. Therefore \(L_j\) are actual AFS.

At a fixed point all nonpole terms have finite limits. If some \(p_i>0\)
and \(f_i=0\), then \(q_{ij}\) is eventually positive and
\(q_{ij}F_{ij}\to-\infty\), while the other terms are bounded above.
Thus the limit is the desired support even on its \(-\infty\) locus.
Omit zero coefficients entirely.

The individual floors are essential. On \(\mathbb A^2_{z,w}\), the
function \(\max(p\log|z|,-j)\) is constantly \(-j\) along \(z=0\), even
as \(|w|\to\infty\). It cannot contain a positive retained
\(w\)-generator branch. Flooring the entire singular sum loses precisely
that growth; (4.6) does not.

\subsubsection{4.4. Countable strict potentials and an explicit
diagonal}\label{countable-strict-potentials-and-an-explicit-diagonal}

Suppose \(k\) is countable, algebraically closed and trivially valued.
Enumerate its affine coordinate algebra and retain embedding generators
in every list. These finite tropicalizations generate the analytic
topology. For an AFS exhaustion \(H\ge1\), degree bounds and (4.2) allow
positive rational \(\gamma_j\) such that
\[0\le\gamma_j\chi_{N_j}\circ\operatorname{Trop}_j\le2^{-j}(H+1).\] The
sum \(\rho=\sum_j\gamma_j\chi_{N_j}\circ\operatorname{Trop}_j\) is
continuous by uniform convergence on each generator-bounded compact, and
has upper logarithmic growth. To prove actual decreasing AFS
regularization, set
\[T_r=\sum_{j\le r}\gamma_j\chi_{N_j}\circ\operatorname{Trop}_j+2^{-r}(H+1).\]
The functions \(T_r\) decrease. Approximate each finite sum uniformly on
the whole affine space by an AFS \(A_r\) with error
\(\epsilon_r=4^{-r}\) using (4.5), and put
\(B_r=2\sum_{i\ge r}\epsilon_i\). Then \(A_r+B_r\) decreases, since its
consecutive difference is at least
\(2\epsilon_r-\epsilon_r-\epsilon_{r+1}>0\). Its limit is \(\rho\).

For \(x\in Y^{\mathrm{an}}\) and a neighborhood \(B\), a finite
tropicalization has an inverse-image neighborhood contained in \(B\).
Replace its one \(\chi\) summand by (4.3). The resulting \(\rho'\)
satisfies \[\rho'\le\rho,\qquad \rho'(x)=\rho(x),\qquad
 \rho'\le\rho-\delta\ \text{outside }B. \tag{4.8}\] It has an actual
decreasing AFS sequence. For \(r\) beyond the replaced index,
approximate the remaining finite continuous summands as above, replace
the fixed support by its \textbf{decreasing} sequence (4.7), retain the
tail \(2^{-r}(H+1)\), and add \(B_r\). Consecutive differences are
nonnegative from the decreasing support and the preceding error
estimate. The other summands stay bounded above at each pole of the
support, proving the correct extended limit. This explicit diagonal
avoids any unsupported uniform approximation of a real support at its
logarithmic poles. The construction is a countable auxiliary theorem,
not an arbitrary-cardinality envelope theorem.

\subsection{5. Original-ring hull localization using actual pole
orders}\label{original-ring-hull-localization-using-actual-pole-orders}

Let \(A\) be an integral affine \(k\)-algebra over the retained
trivially valued \(k\), \(f\ne0\), \(h\in A_f\), and \(U\) an AFS
function on \(Y=\operatorname{Spec}A\). Assume, for \(a,c>0\) and
\(b,C\in\mathbb R\),
\[a\log|h|+b\le\max(U,-c\log|f|)+C\quad\text{on }Y_f^{\mathrm{an}}. \tag{5.1}\]
The finite normalization \(\bar A\) is used only for algebraic pole
orders. For a height-one divisor \(D\) of \(\bar A\), follow the
divisorial ray \(r\operatorname{ord}_D\), \(r\to\infty\). Regular
elements of \(A\) have nonnegative order, so the finitely many branches
of \(U\) stay bounded above there. Applying (5.1) gives
\[a\max(-\operatorname{ord}_Dh,0)\le c\operatorname{ord}_Df. \tag{5.2}\]
Outside \(Z(f)\), \(h\) is regular. For positive integers \(p,q\) with
\(q/p>c/a\), (5.2) shows that \(z=h^pf^q\) has no height-one poles.
Normality identifies \(\bar A\) with the intersection of its height-one
valuation rings, so \(z\in\bar A\) and is integral over \(A\). Since
\(z\in A_f\), the finite module \(A[z]/A\) vanishes after localization
at \(f\). One fixed \(f^N\) kills it. Therefore
\[f^Nz^l=f^{N+ql}h^{pl}\in A\quad(l\ge1). \tag{5.3}\]

Let \(L_0\subset Y_f^{\mathrm{an}}\) be compact and let
\(x\in\widehat L_{0,Y}\cap Y_f^{\mathrm{an}}\), where the hull is
defined using \textbf{the original ring \(A\)}. Suppose
\(|\log|f||\le M\) on \(L_0\cup\{x\}\). Use (5.3) as actual regular hull
tests. When \(h(x)\ne0\) and \(\sup_{L_0}|h|>0\), dividing the log
inequality by \(pl\), then letting \(l\to\infty\), gives
\[a\log|h(x)|+b\le\sup_{L_0}(a\log|h|+b)+2a(q/p)M.\] Let
\(q/p\downarrow c/a\) to obtain
\[a\log|h(x)|+b\le\sup_{L_0}(a\log|h|+b)+2cM. \tag{5.4}\] If \(h(x)=0\)
this is immediate. If \(h\) vanishes on all of \(L_0\), the regular
tests (5.3) force \(h(x)=0\) because \(f(x)\ne0\). No subtraction of
infinities is used.

Denominator exponent counting is invalid. In \(A=k[z]\), take \(f=z^2\),
\(h=z^{-1}=z/f\), \(a=2,c=1\), and \(U=\max(0,2\log|z|)\). Then (5.1)
holds with \(b=C=0\) everywhere on \(Y_f\), but \(a\cdot1\le c\) fails
for the displayed denominator exponent one. Actual orders give equality
\(2\cdot1=1\cdot2\). The fixed annihilator in (5.3), rather than a
normalization hull, is the step that puts the tests back in \(A\).

For a decreasing AFS sequence \(w_i\downarrow v\) on \(Y_f\), suppose
the controlled antecedent gives \(w_i\le2\max(U,-\epsilon\log|f|)+B_i\)
for each \(i\), where \(\epsilon>0\) is arbitrary and \(v\le U+C\).
Apply (5.4) to each branch of \(w_i\) with pole budget \(2\epsilon\).
The floor branch is immediate, so
\[w_i(x)\le\sup_{L_0}w_i+4\epsilon M.\] For continuous functions
decreasing to an upper semicontinuous function on a compact set, their
suprema decrease to the supremum of the limit: nested compact superlevel
sets prove this, including a limit supremum of \(-\infty\). No uniform
Dini assertion is needed. Passing \(i\to\infty\) and then
\(\epsilon\downarrow0\) gives \[v(x)\le\sup_{L_0}v. \tag{5.5}\] Equation
(5.4) is proved independently. Equation (5.5) is conditional on the
precisely stated controlled-approximation antecedent; it is not a
consequence of the printed lower-growth lemma alone.

\subsection{6. The repaired split-piece
implication}\label{the-repaired-split-piece-implication}

This section is a conditional mathematical theorem. Its antecedents are
stated to prevent an unreviewed generated proof from being silently
accepted.

\textbf{Metric antecedent.} An actual AFS reference admits a rational
ample compactification whose affine chart is the original affine variety
and whose local weight is that reference. A continuous rational contact
compact perturbation of such a reference has an actual decreasing AFS
sequence. Continuous glued rational ample metrics with a whole
reference-winning collar are genuinely global PSH; decreasing closure
and already-genuine same-ample-class sequential FS regularization, with
exact section/generator conversion, are available. Controlled
approximation in the precise form used for (5.5) is available. No
upper-envelope property on a singular target is included.

\textbf{Smooth antecedent.} Smooth upper-growth weak affine
regularization and principal-open removal are proved using target-side
smooth EP, including equality of the removed-set values on the original
witness.

\textbf{Trace antecedent.} For finite surjective \(\pi:Y\to S\), with
\(Y\) integral affine and \(S\) a smooth affine open of \(\mathbb A^n\),
there are input-independent positive integer weights \(m_y(s)\) on every
finite analytic fiber, summing to the generic degree \(d\). Every weak
upper-growth input with a single proper witness has an AFS-regularizable
trace \[N_\pi u(s)=\sum_{y\in\pi^{-1}(s)}m_y(s)u(y). \tag{6.1}\] For
genuine AFS tests this is the actual determinant norm extended over the
normal base, hence continuous everywhere. The weights include the
correct slicing and residue contributions; on a nonflat fiber they are
not assumed to equal arbitrary special-fiber algebra dimensions.
Complete-field base change and boundary radial identification are part
of this antecedent.

None of these complete antecedents has been established anew in this
paper. Sections 4--5 and the argument below independently check the
additional implications the proposal needs.

\subsubsection{6.1. Rossi and the localized hull
implication}\label{rossi-and-the-localized-hull-implication}

The metric antecedent proves the continuous Rossi boundary lemma: if
\(x\in\widehat L_Y\setminus L\) and \(B\Subset Y^{\mathrm{an}}\) is an
open neighborhood of \(x\) disjoint from \(L\), then
\[x\in\widehat{\partial B\cap\widehat L_Y}_Y. \tag{6.2}\] To see the
actual compact perturbation, if a polynomial separates \(x\) from the
indicated boundary set, separate the compact remaining bad boundary part
from \(\widehat L_Y\) by finitely many original polynomials. Their
maximum with zero and generator bounds gives an AFS function \(g\) which
vanishes on the hull and is positive on that bad boundary. Multiply
\(g\) by a large integer so it wins over \(\log|h|-t\) on the whole
boundary. It wins on a whole collar by compactness and continuity (the
logarithmic pole causes no problem because the finite branch wins
there). Glue its maximum with \(\log|h|-t\) inside \(B\) and \(g\)
outside. A uniform AFS approximant with error less than a third of the
strict separating value at \(x\) contradicts the original hull
inequalities. Uniform approximation on a compact set follows from the
metric antecedent and continuity of this particular glued function; it
is not a Dini assertion for an upper semicontinuous target.

Let a weak upper-growth \(u\) be already locally AFS-regularizable on
affine opens covering \(Y\setminus Z\). For rational \(c>0\), the
repaired localized implication is
\[ (u+c\rho)(x)\le\sup_L(u+c\rho)\quad(x\in\widehat L_Y\setminus Z^{\mathrm{an}}). \tag{6.3}\]
The hypothesis concerns the original \(u\), not merely \(u+c\rho\). If
(6.3) fails, let \(v=\max_i\log|f_i|\) for generators of the ideal of
\(Z\). Individual decreasing generator floors show that \(v\) is
AFS-regularizable and has \(-\infty\) locus \(Z\). Choose small rational
\(\epsilon>0\) so \(u+c\rho+\epsilon v\) still separates a tested point
from \(L\). It attains a finite maximum on the compact hull at
\(y\notin L\cup Z\). Choose a relatively compact \(B\) about \(y\) with
closure in a principal open of the local regularization chart and
disjoint from \(L\). Replace \(\rho\) by (4.8). The original local
hypothesis regularizes \[q=u+c\rho'+\epsilon v\] on that principal open.
It has global AFS upper growth there, and
\(q(y)>\sup_{\partial B\cap\widehat L_Y}q\) with a drop of \textbf{\(c\)
times} the barrier gap. Apply (6.2) and then (5.5) to the compact
boundary set, which lies in that same principal open. The contradiction
proves (6.3). Subtracting a psh barrier from a regularizable sum would
not justify this step.

\subsubsection{6.2. Localize the integer trace without moving the
original hull
neighborhood}\label{localize-the-integer-trace-without-moving-the-original-hull-neighborhood}

Assume \(\pi|_Z\) is a closed immersion into its scheme image and
\(Z'=\pi^{-1}\pi(Z)\setminus Z\) is closed and disjoint from \(Z\).
Suppose \(\{u+c\rho<0\}\) is not Runge. Its definition gives compact
\(L\) in that set whose compact hull contains a point where
\(u+c\rho\ge0\). Normalize the finite maximum on that hull to zero and
choose \(x\) attaining it. Pick rational \(\delta>0\) with
\(\sup_L(u+c\rho)\le-\delta\). By (6.3), \(x\in Z^{\mathrm{an}}\).

Let \(C=\widehat{\{\pi(x)\}}_S\). Every polynomial vanishing on
\(\pi(Z)\) vanishes on \(C\), so \(C\subset\pi(Z)^{\mathrm{an}}\). Its
preimage splits into compact pieces in \(Z\) and \(Z'\). Choose disjoint
analytic neighborhoods \(A_1,A_2\) of these pieces, with \(A_1\)
avoiding \(Z'\) and \(A_2\) avoiding \(Z\). Proper closedness gives an
initial neighborhood \(A\) of \(C\) with
\(\pi^{-1}A\subset A_1\cup A_2\).

Construct finitely many polynomial constraints with strict margins on
\(C\), including generator bounds, whose compact closed intersection
lies in \(A\). If their log bounds are \(b_j\), choose rational \(e>0\)
below every margin and put
\[g=\max\left(0,\frac{\log|a_j|-b_j+e}{e}\right)_j.\] Then
\(Q=\{g\le1\}\) is compact in \(A\) and \(g=0\) on an open neighborhood
\(U\) of \(C\). Set \(W=\{g<1\}\) and \(W_i=\pi^{-1}W\cap A_i\). These
are open with compact closures, not compact open sets.

On \(W_1\) let \(h=1\), on \(W_2\) let \(h=0\). Glue
\[\widetilde h=\max(3g\circ\pi-1,h)\text{ on }\pi^{-1}W,
 \quad\widetilde h=3g\circ\pi-1\text{ elsewhere}. \tag{6.4}\] For
\(g>2/3\) the outer reference wins on both pieces. This is a whole
collar; the metric antecedent gives a continuous AFS-regularizable
function.

The trace is continuous on the \textbf{whole} base. On a compact base
neighborhood \(E\), its genuine AFS approximants converge uniformly on
compact \(\pi^{-1}E\). For any two bounded inputs the fixed positive
weights give
\[\|N_\pi a-N_\pi b\|_E\le d\|a-b\|_{\pi^{-1}E}. \tag{6.5}\] Each AFS
trace is a continuous extended determinant formula on the normal base,
so their locally uniform limit is continuous. Continuity on an étale
open alone is insufficient.

On \(U\), \(N_\pi\widetilde h\) is the sum of the positive integer
weights in \(\pi^{-1}U\cap A_1\). It is locally constant. Let \(\gamma\)
be its value at \(\pi(x)\) and choose an open \(O\subset U\) about
\(\pi(x)\) on which that value is \(\gamma\). The initial
\(A,W,W_i,Q,g\) are kept fixed.

If \(\gamma>1\), upper semicontinuity gives a neighborhood \(N\) of
\(x\) in \(W_1\) on which \(u+c\rho<\delta/[2(\gamma-1)]\). The
\(W_1\)-fiber at \(\pi(x)\) is just \(x\), since all other points over
\(\pi(Z)\) lie in \(Z'\) and \(A_1\) avoids \(Z'\). In \(\pi^{-1}O\),
the \(W_1\) piece is clopen. Proper closedness therefore gives a smaller
\(O'\subset O\) whose entire \(W_1\) preimage lies in \(N\). If
\(\gamma=1\), set \(O'=O\) and omit this estimate.

Choose open \(B\) about \(x\) with compact closure contained in
\(W_1\cap\pi^{-1}O'\) and disjoint from \(L\). Put
\(K=\partial B\cap\widehat L_Y\). This proves, before any later
estimate, that \textbf{every} point of \(K\), and every \(W_1\) point of
its base fibers, lies in the shrunk region where the same integer
\(\gamma\) applies. The choice shrinks only the test-fiber neighborhood,
not the original polynomial hull neighborhood or the gluing domains.

\subsubsection{6.3. One witness, a winning collar, and both uniform
bounds}\label{one-witness-a-winning-collar-and-both-uniform-bounds}

Use (4.8) at \(x\) in \(B\) and put \(u'=u+c\rho'\). Then \(u'(x)=0\),
\(u'\le u+c\rho\), and \(u'\le-e'\) on \(K\) for some \(e'>0\)
containing the factor \(c\).

Set \(v=u'\) on \(W_1\) and zero on \(W_2\). Choose rational
\(C_0>\sup_{\pi^{-1}Q}\max(u',0)+\delta+1\) and define
\[\widetilde v=\max(C_0g\circ\pi-\delta,v)\text{ on }\pi^{-1}W,
 \quad\widetilde v=C_0g\circ\pi-\delta\text{ elsewhere}. \tag{6.6}\]
Since \(C_0-\delta\) exceeds the bounded upper value of \(v\) on compact
\(\pi^{-1}Q\), the outer reference wins on a whole strip before \(g=1\).
Weak membership uses the \textbf{single original proper witness for
\(u\)}: pull back \(\rho'\) and \(g\), form the maxima on the disjoint
preimage pieces and glue on this common winning collar. No unrelated
witnesses are glued. The resulting function has global upper growth and
is generically finite because of its finite floor. On \(\pi^{-1}O'\),
\[\widetilde v=\max(u',-\delta)\text{ on }W_1,\qquad
 \widetilde v=0\text{ on }W_2.\]

By the trace antecedent, \(w=N_\pi\widetilde v\) is regularizable on
\(S\). Its finite pullback is actual AFS-regularizable on \(Y\). To
retain source generators, use their monic integral equations: if
\(z^m+\sum_{i<m}b_i z^i=0\), the non-Archimedean norm bound is
\[\log|z|\le\max_{i<m}\frac{\log|b_i|}{m-i}.\] Base polynomial degree
bounds and retained base-generator lower growth make a small positive
source-generator branch dominated by each pulled AFS function. Append it
without changing that function. This corrects the denominator in the
printed root bound as well as the generator omission.

Rossi (6.2) gives \(x\in\widehat K_Y\). Actual AFS hull inequalities and
decreasing compact suprema therefore give
\[0=w(\pi(x))\le\sup_K w\circ\pi. \tag{6.7}\] We obtain uniform negative
bounds on all of \(K\). If \(y\in K\cap Z\), the \(W_1\) fiber is the
singleton \(y\) of weight \(\gamma\), so
\[w(\pi(y))\le-\gamma\min(\delta,e').\] If \(y\in K\setminus Z\), (6.3)
for the original \(u+c\rho\) gives \(u'(y)\le-\delta\), so
\(\widetilde v(y)=-\delta\) with weight at least one. If \(\gamma>1\),
every other \(W_1\) point in that fiber has value less than
\(\delta/[2(\gamma-1)]\) and the total remaining weight is at most
\(\gamma-1\). The \(W_2\) contributions vanish. Hence
\[w(\pi(y))\le-\delta/2.\] If \(\gamma=1\), there is no other
positive-weight \(W_1\) point, so the bound is \(-\delta\). These bounds
contradict (6.7). This proves the split-piece Runge conclusion under the
stated antecedents.

Normalize \(\rho\ge0\) and let rational \(c\downarrow0\). The sublevels
\(\{u+c\rho<0\}\) increase to \(\{u<0\}\). Every compact subset of that
union lies in one member; hull containment then proves that the union is
Runge. There is no global integer-constancy assertion on a disconnected
base open.

\subsection{7. Strict Hartogs domains and finite-projection
geometry}\label{strict-hartogs-domains-and-finite-projection-geometry}

\subsubsection{7.1. Runge comes before the radius
formula}\label{runge-comes-before-the-radius-formula}

In the countable-field setting, an affine analytification is second
countable and locally compact, hence has compact exhaustions. If an open
\(D\subset(Y\times\mathbb A^1)^{\mathrm{an}}\) is Runge, the hulls of
its compact exhaustion sets lie compactly in \(D\). On a larger
generator box, every point outside a compact hull can be separated by a
polynomial with a strict norm gap. Compactness of the bad complement
gives finitely many such polynomials. Rational thresholds inside the
gaps and additional generator bounds give a strict polynomial
neighborhood with compact closed bounds contained in \(D\). Repeating
with each prior closed bound included in the next compact exhaustion set
yields \[Q_i=\{|g_{ij}(y,w)|<e^{b_{ij}}\ \forall j\},\qquad
 \overline Q_i\subset Q_{i+1},\qquad\bigcup_iQ_i=D, \tag{7.1}\] where
the notation \(\overline Q_i\) here can be replaced by the explicitly
constructed compact weak-inequality bound; no equality with that bound
is required. Include original base-generator bounds and a bound on
\(w\).

Suppose \(g_{ij}=\sum_l f_{ijl}w^l\). At a point where \((y,0)\in Q_i\),
the constant coefficients satisfy their strict inequalities. The largest
fiber disk radius \(d_{Q_i}(y)\) has the exact formula
\[-\log d_{Q_i}(y)=\max_{j,l\ge1}\frac{\log|f_{ijl}(y)|-b_{ij}}{l}. \tag{7.2}\]
Indeed every point of a disk satisfies the weighted Gauss bound.
Conversely Gauss seminorms at arbitrary positive real radii exist over
every complete residue field, including a trivially valued one, and test
each coefficient. The supremal open radius is therefore exactly (7.2).
The explicit \(w\) bound supplies a finite floor. On a fixed
generator-bounded box within the zero-section domain, dominated
generator branches retain a genuine local AFS expression. The literal
formula decreases with \(i\).

If \(D\) contains the zero section, its zero-section domains exhaust
\(Y^{\mathrm{an}}\). Every fiber closed disk of radius \(r<d_D(y)\) is
compact in \(D\), so lies in one \(Q_i\). Thus
\(d_{Q_i}(y)\uparrow d_D(y)\) and (7.2) decreases to \(-\log d_D\),
including value \(-\infty\) when the radius is infinite. This argument
assumes and uses Runge before computing coefficients.

The identity between topological interiors of closed weak polytopes and
strict inequalities is false. The constant constraint \(|1|\le1\) has
the entire space as interior, whereas \(|1|<1\) is empty. Strict margins
and (7.1) avoid that error.

For a nontrivial radius potential \(u\le H\), \(H\ge1\) an AFS
exhaustion, the metric antecedent globalizes it without singular EP. On
the projective ample-reference compactification of \(H\), fix a rational
floor \(M\) and use \(\max(u-2H,M)\) with constant extension at the
boundary. Its nonfloor region lies in a common compact box because
\(u-2H\le-H\). Take a larger compact box where the target is below the
floor with a strict margin on the boundary collar. The decreasing
continuous radius expressions (7.2) eventually lie below that floor on
the whole collar: nested compact superlevel sets prove this. They then
glue as genuine continuous metrics in the same ample class. Decreasing
closure gives the floor potential, and then the floor tends to
\(-\infty\). Only after this membership is established may same-class
sequential regularization and section/generator conversion be applied.
This conversion remains conditional on the metric antecedent, not on an
asserted singular upper-envelope theorem.

For \(D=\{\log|w|+u(y)<0\}\) the exact radius is \(e^{-u(y)}\). A
product proper witness gives weak membership. On a chart with supplied
\(u_j\downarrow u\), take individual floors of \(\log|w|\) with a
product exhaustion as in (4.6); positive sums retain all product
generators and give actual local regularization of \(u+\log|w|\). Apply
the split-piece Runge theorem on the product \textbf{first}, then (7.2)
and the preceding collar conversion.

\subsubsection{7.2. The finite projection at all residue
conjugates}\label{the-finite-projection-at-all-residue-conjugates}

The following algebraic lemma is proved independently. Let \(K\) be an
infinite characteristic-zero field, \(Y\) integral affine of dimension
\(n\), and \(Z\subsetneq Y\) reduced closed and smooth at a closed point
\(z\), with \(\dim Z\le n-1\). There is an affine neighborhood \(U\) of
\(z\), an affine open \(S\subset\mathbb A_K^n\), and a finite surjective
\(\pi:U\to S\) whose restriction to \(Z\cap U\) is a closed immersion
into its scheme image.

Embed a projective compactification \(X\) in \(\mathbb P^d\), \(d>n\),
and let \(T\) be the reduced union of the closure of \(Z\) and
\(X\setminus Y\). It is smooth at \(z\). Over an algebraic closure,
there are finitely many residue conjugates \(z_\alpha\). For every
\(z_\alpha\), avoid the union of lines joining \(z_\alpha\) to other
points of \(T\), its closure, and the projective tangent directions at
\(z_\alpha\). Each incidence locus has dimension at most
\(\dim T+1\le n<d\); include lines between distinct conjugates. Avoid
also \(X\), of dimension \(n\). Their finite union is a proper closed
set. It is invariant under the finite separable conjugacy action and
descends to a proper \(K\)-closed set. Since \(K\) is infinite, a
\(K\)-rational center exists outside it.

Projection from that center is finite on \(X\): a positive-dimensional
fiber would contain a projective line through the center lying in \(X\),
contradicting its exclusion. On \(T\) its geometric fiber at each
conjugate image is one point and its tangent map is injective. Thus the
local geometric fiber algebra is the residue field: the kernel of the
tangent map gives \(\mathfrak m_B=\mathfrak m_A B+\mathfrak m_B^2\), and
local Nakayama gives \(\mathfrak m_B=\mathfrak m_A B\). Distinct
conjugate images guarantee no residue-field collapse. Equivalently one
may do this entire fiber calculation after faithfully flat extension to
the algebraic closure.

Let \(A\to B\) be the finite local algebra of the restriction to its
\textbf{scheme image}, so the image ring injects into \(B\). The fiber
map is surjective onto the residue field; hence the finite cokernel
satisfies \((B/A)\otimes\kappa(\mathfrak m_A)=0\). Nakayama gives
\((B/A)_{\mathfrak m_A}=0\). On a neighborhood of that image point the
finite map is an isomorphism. This is a cokernel argument in the image
ring, not a bare tangent-space claim.

Repeat the construction down to \(\mathbb P^n\), preserving the local
isomorphism at the conjugates by shrinking the image neighborhoods. The
composite is finite surjective onto \(\mathbb P^n\): its closed image
has dimension \(n\). The point's fiber on \(T\) contains no boundary
point. Remove the closed image of \(X\setminus Y\) and the bad
restriction locus, and choose a principal affine neighborhood \(S\) of
the image inside a standard \(\mathbb A^n\) chart. Set
\textbf{\(U=\pi^{-1}S\)}, the actual preimage of this chosen target.
Finiteness makes \(U\) affine and \(U\subset Y\). This proves the
claimed source and target restrictions. It also corrects the
larger-preimage mistake in Wu LF2656--2658.

\subsubsection{7.3. The two terminating inductions, conditional on the
analytic
gates}\label{the-two-terminating-inductions-conditional-on-the-analytic-gates}

For a fixed finite projection injective on \(Z\), define reduced closed
sets
\[Z_1=Z,\qquad Z_{j+1}=Z_j\cap\overline{\pi^{-1}\pi(Z_j)\setminus Z_j}^{\,\mathrm{Zar}}. \tag{7.3}\]
Finiteness preserves dimension. Every maximal-dimensional component of
\(Z_j\) is a component of \(\pi^{-1}\pi(Z_j)\) and is not contained in
the closure of the complementary components. Thus
\(\dim Z_{j+1}<\dim Z_j\) whenever \(Z_j\ne\varnothing\). On the
complement of \(\pi^{-1}\pi(Z_{j+1})\), the pieces over \(Z_j\) are
closed and disjoint. Apply the split theorem to \(u+\log|w|\) on the
actual affine preimage of a principal base neighborhood there, then use
the Hartogs conversion. Points over \(\pi(Z_{j+1})\) outside \(Z_{j+1}\)
are outside the original \(Z\) by injectivity on \(Z\) and retain the
original local regularization hypothesis. This proves the
descending-\(Z\) induction under the stated analytic gates. The
conclusion is local on those preimages, not global on an unlocalized
ambient variety.

Next take successive \textbf{reduced} singular loci of \(Y\). In
characteristic zero each next singular locus has smaller dimension, so
this finite filtration terminates. The smooth open stratum uses the
smooth antecedent. At a point of the next smooth stratum, the
finite-projection lemma and (7.3) supply local regularization from the
preceding stratum's affine cover. The cover can be formed from closed
points because the complement of the union of their Zariski
neighborhoods, if nonempty, has a closed point. Principal shrinking and
Section 3 retain the affine hypotheses. Finally, with \(Z=\varnothing\),
(6.3) proves the global sublevel Runge property; applying it to the
product Hartogs function and then (7.2) proves global AFS
regularization. The final empty-\(Z\) step is essential and is now
explicit.

On projective ample charts retain the actual reference weights
\(\rho_i\). For a very ample standard chart system,
\(\max_i(-\rho_i)=0\). Damp a weak local weight by \((c-1)\rho_i\), glue
its boundary \(-\infty\) potential using the preceding affine conversion
and whole-floor collars, and take the finite maximum of the genuine
potentials in class \(c\omega\). This equals the original weak
potential. Let rational \(c\downarrow1\) using genuine class closure.
This gives countable weak/genuine equality \textbf{if} the analytic
gates are proved, not an unconditional theorem in this deliverable.

\subsection{8. Cardinality, normalization and the remaining order of
proof}\label{cardinality-normalization-and-the-remaining-order-of-proof}

The preceding conditional chain is not the original arbitrary-\(k\)
theorem. A countable-field envelope base must first be proved. Even for
countable \(k\), one needs constancy on every geometric fiber of a
resolution for a genuine psh potential in the pulled-back nef class, and
upper semicontinuous descent equal to the divisorial minimum extension.
The printed Wu 7.1 proof at LF2705--2717 invokes both Proposition 2.24
and weak/global 6.5. This paper does not assume them.

There is a precise conditional cardinality reduction, not a descent of
every competitor to one field. Suppose countable rational ample EP has
been proved for the original geometrically unibranch domain, with
genuine continuous-envelope orthogonality, positive MA base change, and
domination. For a finite algebraic PL obstacle \(f\) and ample rational
reference, descend its finite data to a countable algebraically closed
\(k_1\subset k\). Ensure the geometric hypotheses by also descending
finite normalization data and a dense-open birational witness; normality
and a finite universal homeomorphism descend faithfully flatly. Let
\(p_1=P(f_1)\) be the countable continuous envelope, and pull it to
\(p\) over \(k\).

Any single model psh competitor \(\psi\le f\) descends, together with
its finite nef model data, to a larger countable algebraically closed
\(k_2\subset k\). The direction is \(\psi=\psi_2\circ r_{k/k_2}\), not
\(\psi_2=\psi\circ r\). Analytic projection is surjective, so the
obstacle inequality descends. Positivity of MA and exact pushforward
imply \[\int(f_2-p_2)\,dM^{\mathrm{raw}}(p_2)
 =\int(f_1-p_1)\,dM^{\mathrm{raw}}(p_1)=0.\] The integrand is
nonnegative continuous, so the genuine measure is supported on its
contact set. On that support \(\psi_2\le f_2=p_2\); genuine domination
gives \(\psi_2\le p_2\). All finite model competitors are therefore
below \(p\), and their supremum gives the large-field envelope. General
decreasing competitors need not descend together. Rational ample EP then
passes to a real ample class via rational classes approaching it from
above, exactly BJ 5.9 and class closure.

The base-change antecedent must use actual relative closure. For a
relatively integrally closed test model over an algebraically closed
smaller field, central integral components stay geometrically integral
on field extension, their generic vertical DVRs have unchanged
\(\operatorname{ord}_E(t)\), and intersection degrees are unchanged. To
justify that the base-changed model remains relatively integrally
closed, one must retain vertical \(R_1,S_2\) under the flat
geometrically regular scalar extension, then apply BJ 1.13; the generic
fiber is the original base-changed nonnormal variety. Model formula
(2.1) then pushes forward with exact \textbf{positive} mass. Uniform
continuous approximation and finite-energy decreasing approximation pass
it to the genuine operator. An absolute normalization or a signed
pushforward without positivity would not justify contact support. This
is the necessary algebraic-and-analytic base-change route; its complete
infinite-extension and descent audit is not claimed here.

For geometrically unibranch integral \(X\) in characteristic zero,
normalization \(q:X^\nu\to X\) has singleton geometric fibers and purely
inseparable residue extensions, hence unchanged residue fields. It is a
finite universal homeomorphism. The norm extension is unique on each
residue field, so analytification is bijective; properness on compact
Hausdorff spaces gives a homeomorphism. Unconditionally for genuine
downstairs inputs, BJ pairing and model projection give
\[E_{q^*\omega}(q^*\phi)=E_\omega(\phi),\qquad
 q_*M^{\mathrm{raw}}_{q^*\omega}(q^*\phi)=M^{\mathrm{raw}}_\omega(\phi). \tag{8.1}\]
Onto PSH is separate. BJ 5.13 requires approaching \textbf{target-side}
larger classes with EP. Once EP is proved on \(X\), take
\((1+1/m)\omega\) to meet this hypothesis. Only then do onto
\(\mathcal E^1\) and equality of finite-dual-energy target sets follow.
Topological normalization cannot substitute for that step.

After actual arbitrary-cardinality EP, Section 2 supplies all retained
compactness, energy differentiation, comparison, attainment and
uniqueness inputs for the original arbitrary-support equation. At
present its geometric antecedent is unproved. A countable rational local
theory, an auxiliary inverse measure or a formal completion does not
close the equation.

\subsection{9. Diagnostics and exact unresolved
assertions}\label{diagnostics-and-exact-unresolved-assertions}

The printed lower-growth Proposition 4.19 is false. On
\(\mathbb A^1_k\), let \(r=\max(\log|z|,0)\) and \(v=r^2\). It is
continuous and locally psh even under the corrected rational contact
definition: on \(|z|<e^M\), rational \(M>0\), interpolate \(t^2\) by
rational secants on nested dyadic partitions of \([-M,0]\), extend by
zero for \(t\ge0\), and put \(t=-\log|z|\). Convexity of the quadratic
makes the secants decrease under refinement. Their finite affine
branches and the zero branch are rational; near \(t=\infty\) they are
constant. Thus they form the required local decreasing sequence. Since
\(r\le r^2+1/4\), the actual AFS \(r\) satisfies Wu's \textbf{printed
lower-growth} condition LF766--768. On the Gauss ray of radius \(e^s\),
\(s\ge0\), every polynomial log norm equals its degree times \(s\).
Every global AFS is bounded by \(As+B\), whereas \(v=s^2\). A decreasing
AFS sequence for \(v\) would have a first term majorizing the quadratic
on the entire ray, impossible. This disproves the printed proposition,
not a corrected upper-growth theorem or the final singular EP assertion.

The repaired constructions in Sections 4--7 remove concrete errors:
missing contact constants, infinite-coordinate supports, real
coefficient poles, denominator counting, omitted scalar \(c\), local
integer weights, division by \(\gamma-1\) at \(\gamma=1\), continuity
only on an étale open, compact-open claims, the wrong winning-locus
formula, witness gluing, topological-interior identities, finite
projection at residue conjugates, the local Nakayama step and actual
target preimages. They establish additional conditional implications;
they do not validate their preceding analytic gates.

The remaining indispensable assertions are:

\begin{enumerate}
\def\labelenumi{\arabic{enumi}.}
\tightlist
\item
  A complete continuous rational metric bridge on the original nonnormal
  projective variety: strict Laurent admissible representation,
  analytic-unit curve sign tests, local-to-global nefness on an already
  represented bundle, algebraization, BE same-class FS approximation,
  exact Gauss coefficient cutoff and compact PL extension. The relevant
  sources allow these applications only after actual PL/model
  representation and gluing. GR Remark 1.8 explicitly does not give an
  arbitrary positive-current pullback or converse.
\item
  Common-collar bounded local products and local comparison/IBP,
  followed by smooth local/global equality, resolved-witness smooth
  upper-growth regularization and exact principal-open removal.
  Already-genuine BJ 12.18 does not establish the initial membership of
  a weak input. The pre-EP controlled approximation used in (5.5)
  remains part of this gate.
\item
  The genuine finite trace antecedent on the whole smooth base,
  including boundary radial identification over every completed residue
  field and input-independent positive integer slicing weights. Finite
  determinant algebra on a locally free open by itself proves none of
  those boundary assertions.
\item
  Resolved-nef fiber constancy and divisorial usc descent, countable
  geometrically unibranch rational EP, and the complete
  arbitrary-cardinality positive base-change/descent argument described
  in Section 8. Normalization onto, unbounded singular envelope
  differentiation and unconditional variational attainment consequently
  remain conditional.
\end{enumerate}

These are uncompleted reconstructions, not nonexistence theorems. The
strongest unconditional new mathematical drafts here are the explicit
barrier/real-support and original-ring localization theorems, the strict
coefficient radius formula on an already Runge domain, and the
finite-projection lemma. The original existence theorem remains
unresolved. Candidate=false.

\subsection{10. Reproducibility and
evidence}\label{reproducibility-and-evidence}

Run \texttt{python3\ rerun.py} in this workspace. The script verifies
the preserved input, checks exact supplied primary hashes and inclusive
one-based LF excerpts, runs bounded finite diagnostics, uses existing
Pandoc to generate standalone TeX, and invokes existing pdfLaTeX twice.
The commands, versions, return codes, stdout/stderr and both pass logs
are retained. No software installation, network request, other model,
worker, campaign/private read, external message, source edit or runtime
change is used. No resource-heavy computation was launched.

\texttt{status.json} has exactly \texttt{id}, \texttt{status},
\texttt{exact\_result}, \texttt{proof\_file}, \texttt{concrete\_gaps}.
\texttt{canonical-notes.json} contains at most four backward-supported
local notes. \texttt{source-extracts/source-provenance.json} contains
literal statements and pertinent proof bodies, exact versions and
hashes, coordinates defined by UTF-8 bytes decoding and
\texttt{split(\textquotesingle{}\textbackslash{}n\textquotesingle{})},
and explicit scope limitations. The whole ten-source text copies are
retained; retention does not claim a whole-source independent audit. The
two standalone paper artifacts are \texttt{P16.1-paper.tex} and
\texttt{P16.1-paper.pdf}.

The finite diagnostics check multinomial identities and maximum bounds
on specified rational grids, exact monotonicity inequalities for
real-support interval approximations, positive integer split-weight
estimates including \(\gamma=1\), and dyadic secant inequalities for the
lower-growth counterexample. These are exhaustive only on their stated
finite grids. Hashes, preservation and compilation certify neither
mathematical correctness nor novelty.

\subsection{References and literal-source
bindings}\label{references-and-literal-source-bindings}

BJ: Boucksom--Jonsson, \emph{Global pluripotential theory over a
trivially valued field}, arXiv:1801.08229v4.

PDF SHA256
\texttt{8d5975337f4ac88a475929ea6a436ec4a15d08715d5582bc0ba0eb3e07152c52}.

Text SHA256
\texttt{659bb8d2e5335781041f03cdb168bac70126487bc0db394f3dac8cf1e9f076d0}.

Wu: Lyuhui Wu, \emph{On the local theory of non-archimedean psh
functions}, arXiv:2609.27610v1.

PDF SHA256
\texttt{11808a6c9ede9c0158b27e16700fc180a4ac8df12a4586a7178e1958e8f6bba8}.

Text SHA256
\texttt{c8ffe91a631da67f6d3872a1a06ab8d15755aed2f926cdf326725ca3662cef7d}.

LF645--707 and 762--774 supply diagnostic definitions; LF950--1038 the
controlled-growth bodies; LF1731--1890 Hartogs/localization;
LF1963--2110 strict barriers/local hulls; LF2200--2372 traces;
LF2373--2533 the entire split proof; LF2534--2692 both inductions;
LF2697--2747 the entire envelope/cardinality proof. These source claims
remain diagnostic where antecedents are not reconstructed.

BE: Boucksom--Eriksson, \emph{Spaces of norms, determinant of cohomology
and Fekete points in non-Archimedean geometry}, arXiv:1805.01016v3.
LF1989--2043 gives exact model/FS differences and common auxiliary
cancellation; LF2334--2352 gives ample nontrivial-field nef/psh
equivalence; LF3125--3177 the continuous Deligne comparison; LF3430
onward and LF3590 onward distinguish bounded relative volumes from
continuous-psh energy. None is an unconditional singular EP theorem.

GM: Gubler--Martin, \emph{On Zhang's semipositive metrics},
arXiv:1608.08030v3. GR: Gubler--Rabinoff, \emph{On semipositive
piecewise linear functions in non-archimedean analytic geometry},
arXiv:2503.05513v1. The relevant GR curve criterion is
\textbf{Proposition} 5.1, Theorem 5.3 and Corollary 5.4, LF818--923;
Remark 1.8 is LF162--164. Actual PL representation is required first.

CLD: Chambert-Loir--Ducros, arXiv:1204.6277v2. Wanner:
arXiv:1801.04713v3. Thuillier: tel-00010990v1. Ducros: published
\emph{Confluentes Mathematici} 4(4), 1250007,
DOI10.1142/S1793744212500077. GJR: Gubler--Jell--Rabinoff,
arXiv:2111.05741v3. Exact hashes and diagnostic proof units for these
bridge sources are in the provenance file; this paper does not claim to
have independently finished their combined new application.

\end{document}
